% Generated by make_tables_v328.py -- do not edit by hand.
\begin{table}
\centering
\scriptsize
\setlength{\tabcolsep}{4pt}
\caption{Selection function of the searched sample against a volume-limited reference census. Reference census: Gaia~DR3 sources with $\varpi>0$, $\sigma_\varpi/\varpi\le0.1$ and $1/\varpi\le40$\,pc (17566 objects; 8594, 49 per cent, carry a \texttt{teff\_gspphot} estimate). Sample: the 88 stars with at least one searched window, at the distances the search itself used. Temperatures are available for 52 of the 88 stars (all 88 matched to the 168-entry ALMA-covered census, 58 by exact name; 52 of those matches carry a temperature), so the temperature panel is normalised within each column over classified objects only, with $T_{\rm eff}$ a proxy for spectral class. ``Ratio'' is the sample column fraction over the census column fraction: 1 means the sample tracks the reference population, $>1$ over-representation.}
\label{tab:selfunc}
\begin{tabular}{@{}lrrr@{}}
\toprule
Property & Census & Searched & Ratio \\
\midrule
\multicolumn{4}{@{}l}{Distance (all 88 searched stars; 17566 census objects)} \\
\quad $d\le5$\,pc & 48 (0.3\%) & 12 (13.6\%) & 49.9 \\
\quad 5--10\,pc & 240 (1.4\%) & 13 (14.8\%) & 10.8 \\
\quad 10--20\,pc & 2069 (11.8\%) & 18 (20.5\%) & 1.7 \\
\quad 20--30\,pc & 5391 (30.7\%) & 18 (20.5\%) & 0.7 \\
\quad 30--40\,pc & 9818 (55.9\%) & 27 (30.7\%) & 0.5 \\
\midrule
\multicolumn{4}{@{}l}{$T_{\rm eff}$ class$^{a}$ (52 searched; 8594 census)} \\
\quad A ($\ge7500$\,K) & 24 (0.3\%) & 1 (1.9\%) & 6.9 \\
\quad F (6000--7500\,K) & 402 (4.7\%) & 9 (17.3\%) & 3.7 \\
\quad G (5300--6000\,K) & 860 (10.0\%) & 16 (30.8\%) & 3.1 \\
\quad K (3900--5300\,K) & 1400 (16.3\%) & 5 (9.6\%) & 0.6 \\
\quad M ($<3900$\,K) & 5908 (68.7\%) & 21 (40.4\%) & 0.6 \\
\midrule
\multicolumn{4}{@{}l}{Other axes} \\
\quad Known planet hosts$^{b}$ & 20 (11.9\%) & 18 (20.5\%) & 1.7 \\
\quad Searched stars / systems & -- & 88 / 81 & -- \\
\bottomrule
\end{tabular}

\smallskip
{\footnotesize $^{a}$Bins are temperature bins, not MK types: A $\ge7500$, F 6000--7500, G 5300--6000, K 3900--5300, M $<3900$\,K. Gaia \texttt{teff\_gspphot} is unavailable for about half the reference census, preferentially for the coolest and faintest objects, so the census M fraction is a lower limit and the earlier-type ratios are correspondingly upper limits. White dwarfs are not separable in this parameterisation and fall wherever their (unreliable) photometric temperature places them.\\ $^{b}$The Gaia reference census carries no planet information; the census column here is instead the 168-entry ALMA-covered census, of which 20 entries are known planet hosts. The 88 searched stars include 18 hosts of 41 known planets.}
\end{table}
